\chapter{Architecture Overview}
\ChapterLabel{ArchitectureOverview}
\section{Chapter Responsibility}
\label{Section:ArchitectureOverviewResponsibility}
Explain the major design choices and why they address the design drivers.

\noindent\textbf{Expected result:} A defensible design direction with explicit trade-offs.

\section{Revision 0.2 Scope Notice}
\label{Section:DeferredArchitectureOverview}
\begin{fsaNote}[Revision 0.2 project content]
\ArtifactRef{Proposal}{REQ-OMS-01} records an open architecture proposal.
The Decision Template and Example section provides template guidance.
Detailed architecture views are deferred under
\hyperref[Section:ProjectScope]{the revision scope}. The
\hyperref[Table:ProjectGapReview]{reviewed gap list} records project decisions
and unresolved items.
\end{fsaNote}

\section{OMS Transfer Proposal (REQ-OMS-01)}
\label{Proposal:REQ-OMS-01}
\index{REQ-OMS-01}
\index{Open Mission Systems}
\noindent\textbf{Decision pending.} Joel Murphy and Jason Lefler own the
adoption decision. Consider \gls{fsaoms} if the
\hyperref[Dependency:Consumer]{downstream processing consumer} needs standard
product discovery. The scope is the host-to-consumer interface in
\hyperref[Table:ProjectInterfaces]{Initial Interface Context}.

\paragraph{Proposed version baseline.}
Evaluate OMS 2.5 (OMSC-STD-001 Rev M, January 22, 2026) with \gls{fsauci}
2.5 and its 2.5.0 schema (OAC-STD-002 Rev E), the schema baseline cited by OMS
2.5. Treat UCI 2.6 or a later release as a separate compatibility decision.
Pin the exact schema artifact and CAL implementation release in the interface
contract if the proposal is adopted \cite{OMSStandard25,UCIRelease26}.

\paragraph{CAL profile.}
Evaluate the Language-Agnostic Critical Abstraction Layer (\acrshort{fsacal})
under OMSC-SPC-013 Rev B.
Use WebSocket with \gls{fsaowp} 1.0 and OMS JSON for UCI publish/subscribe.
Both endpoints would agree on the CAL endpoint, supported schema and topic
names. This profile allows the host and consumer to use different implementation
languages \cite{OMSVDD25,OMSPDD25}.

\paragraph{Proposed UCI messages.}
The host would publish \texttt{ProductMetadata} to describe the acquisition
product and \texttt{ProductLocation} to advertise the corresponding transfer
URI and availability. The downstream processing consumer would subscribe to
those messages and use the URI to open the separate byte transfer.
\texttt{ServiceStatus} is optional for this transfer and could report health
if acquisition is modeled as an OMS Service \cite{OMSStandard25}.

\paragraph{Byte transfer and decision criteria.}
The separate transfer would deliver the same unsigned 8-bit samples unchanged
and in order. The \hyperref[Section:ProjectAcquisitionConditions]{agreed
acquisition conditions} define delivery at consumer interface acceptance and
preservation of bytes successfully read into host memory before a device
disconnection. CAL publication and transport completion do not establish
consumer acceptance. If the endpoints are networked, an ordered TCP stream
is a candidate. No listed OMS data format was identified for the bare ADC0804
stream. The Opaque Data entry is limited to exchanges involving an Opaque
Capability. Evaluate the raw format against the Non-OMS Data Transfer
provisions before claiming a compliance tier \cite{OMSStandard25}.

Compare the OMS option with a direct project-defined consumer interface.
OMS offers standard product discovery at the cost of operating a CAL endpoint
and managing message-schema compatibility. The owners' decision depends on
the consumer's interoperability need and supported CAL profile, acceptance of
the raw format, and a defined consumer acceptance signal, buffering,
backpressure, and reconnection behavior. Evaluate both options against
\RequirementRef{REQ-D-002} and \UseCaseRef{UC-001}.

\section{How to Use This Chapter}
Explain the forces that shape the design, the major decisions, alternatives considered, and consequences. Use the 4+1 views only when they help readers answer different architectural questions.
\section{Decision Template and Example}
\textbf{ARC-001:} \textit{[Decision]}; \textbf{Drivers:} \textit{[quality attributes and constraints]}; \textbf{Rationale:} \textit{[why selected]}; \textbf{Consequences:} \textit{[benefits, trade-offs, and risks]}. Example: use an asynchronous message queue to protect request processing from transient downstream outages; this improves resilience but introduces eventual consistency.
